\section{Combat}
Combat is broken up into rounds during which each participant has the chance to act.
During each round, the GM and Players alternate taking turns to have one creature act, until all participants have acted; if there are more participants on one side or the other the remaining creatures generally act at the end, though one side may pass to allow the other to act an additional time.
Either side may also spend \Advantage{} once each turn to have two creatures go instead of just one.
No creature can act more than once per round, except as allowed by an \Ability{} or other rule.

On a creature's turn, they can move to any point the GM deems they could reasonably reach in a few moments without needing to make a roll and take \textbf{Actions} as described in Section~\ref{sec:actions}.

\subsection{Actions}\label{sec:actions}\index{Action}
There are three kinds of actions that characters can take during combat: \textbf{Main Actions}, \textbf{Quick Actions}, and \textbf{Defend Actions}.
Main Actions and Quick Actions are taken on a character's turn in combat while Defend Actions can be taken when they are attacked.
Each turn, every character gains a certain number of Main and Defend Actions they can use that round (generally one of each) and any unused actions at the end of round are wasted.
Characters can use any number of Quick Actions per turn, though, with the GM's permission.

\paragraph{Main Actions}\index{Main Action|see(Action)}\index{Action!Main}
Main Actions represent significant things that a character can do on their turn.
This includes making an attack, taking more movement or movement that would require a roll (like moving over very rough terrain), \emph{Helping} another character, interacting with a complex object, or other similar actions that the GM may permit.

Additionally, instead of taking a main action, a character can spend a Main Action in the following (or as allowed by another rule):
\begin{itemize}
  \item[-] When making an attack, a character can spend any number of additional Main Actions and to defend against the attack the target must use the same number of Defend Actions.
  \item[-] When making an attack, a character can spend an additional Main Action to target a group of opponents that are clumped together within melee range (at the attacker's size) of each other.
  \item[-] To take a Defend Action.
  \item[-] To gain 2 Advantage.
\end{itemize}
For uses of a main action to boost an attack, the character must declare that use before any rolls for resolving the attack.

\paragraph{Quick Actions}\index{Quick Action|see(Action)}\index{Action!Quick}
Quick Actions represent simple tasks like drinking a potion or drawing a weapon that a character can perform on their turn easily.
A character can take any reasonable number of Quick Actions on their turn, as determined by the GM.

\paragraph{Defend Actions}\index{Defend Action|see(Action)}\index{Action!Defend}
Defend Actions represent a character's ability to defend themself from an attack.
Using a Defend Action impacts the chances of the attack succeed and allows them, if they succesfully avoid the attack, to hit back.
Using a Defend Action also allows the character to activate a Defensive Combat Style which can provide additional bonuses like to the roll or other statistics.

\subsection{Attacks}
The mechanics of an attack depends on whether a Player Character (PC) is attacking a Non-Player Character (NPC) or vice-versa\footnotemark.
\footnotetext{These mechanics do not provide a mechanism for Player vs Player combat. If PvP arises, first the group should discuss if everyone is okay with it and if so what the goal or outcome of the combat will be. Then, these rules can be applied but using Opposed Rolls between the attacking and defending characters.}
The process of resolving an attack in each of these circumstances are described below.

\paragraph{PC Attacking NPC}
\begin{enumerate}
  \item The PC spends a Main Action and identifies the NPC (or NPCs) they are targeting, the weapon and the supported Offensive Combat Style they are using, and whether they are using any special features (such as expending additional Main Actions).
  \item The GM determines, for each NPC being targeted, whether they will expend a Defend Action and if so the Defensive Combat Style they are using. The NPC will have a different difficulty to hit based on whether or not they use a Defend Action (generally the difficulty to hit is $-1$ if the NPC did not use a Defend Action).
  \item The PC then makes a single roll against the difficulty to hit the NPCs, and the attack is resolved for each NPC based on the Outcome:
    \begin{itemize}
      \item \textbf{Very Bad}: The GM gains an Advantage. If the NPC took a Defend Action, and have an attack that could target the PC, they may deal damage to the PC and otherwise the GM gains another Advantage.
      \item \textbf{Bad}: If the NPC took a Defend Action, and have an attack that could target the PC, they may deal damage to the PC and otherwise the GM gains an Advantage.
      \item \textbf{Mixed}: The PC deals damage to the NPC. If the NPC took a Defend Action, and have an attack that could target the PC, they may deal damage to the PC and otherwise the GM gains an Advantage.
      \item \textbf{Good}: The PC deals damage to the NPC.
      \item \textbf{Very Good}: The PC gains an Advantage and deals damage to the NPC.
    \end{itemize}
    Once the PC rolls and determines the outcome, they can choose to spend additional Advantage to gain additional damage dice; if multiple NPCs were targeted by the attack the PC must spend an Advantage for an additional damage die for \emph{each} NPC they wish to do the additional damage to.
\end{enumerate}

When a PC deals damage to an NPC, the damage is determined by their Offensive Combat Style and their Weapon.
A PC can also spend Advantage (including the Advantage they receive on a Very Good result) to add additional damage dice, as described in Section~\ref{sec:advantage}.

\paragraph{NPC Attacking PC}
\begin{enumerate}
  \item The NPC spends a Main Action and identifies the PC (or PCs) they are targeting, the attack they are using (and any of its special features), and whether they are using any special features (such as expending additional Main Actions).
  \item Each PC being targeted determines whether they will expend a Defend Action and if so the Defensive Combat Style they are using.
    \begin{enumerate}
      \item If a PC does not use a Defend Action, the NPC automatically hits (unless the PC has the Armored Defensive Combat Style)\footnotemark and the NPC deals damage to them.
        \footnotetext{This is an asymmetry to the system, PCs always have to roll to hit while NPCs do not. If a group does not like this rule you can instead simply use a higher Difficulty to avoid the hit if the PC does not use a Defend Action. This asymmetric rule is intentional and intended to reflect the danger and chaos of the world and to express that it is not necessarily fair to the PCs.}
      \item If a PC does use a Defend Action they determine the Defensive Combat Style they are using and any associated weapon they are using.
        A PC may also have multiple Defensive Combat Styles that could be helpful, they may spend an Advantage to receive the benefits (other than benefits to the roll) of an additional Defensive Combat Style.
        The PC then gets to roll to avoid the hit using the Basic Statistic of their selected Defensive Combat Style and any bonuses it or their weapon grant, the difficulty of this roll is determined by the NPC's statistics.
        The attack is resolved based on the Outcome:
        \begin{itemize}
          \item \textbf{Very Bad}: The GM gains an Advantage and the NPC deals damage to the PC.
          \item \textbf{Bad}: The NPC deals damage to the PC.
          \item \textbf{Mixed}: The NPC deals damage to the PC. If the PC took a Defend Action, and have an attack that could target the NPC, they may deal damage to the NPC or gain an Advantage.
          \item \textbf{Good}: If the PC took a Defend Action, and have an attack that could target the NPC, they may deal damage to the NPC or gain an Advantage.
          \item \textbf{Very Good}: If the PC took a Defend Action, they gain an Advantage and, if they have an attack that could target the NPC, they may deal damage to the NPC or gain an additional Advantage.
        \end{itemize}
    \end{enumerate}
\end{enumerate}

When an NPC deals damage to a PC, the amount of the damage is determined by their statistics.
The GM can also spend Advantage to add additional damage dice but cannot spend more Advantage for additional damage on a single roll than the current round number (where the first round of combat is round 1).

\paragraph{Taking Damage}
Once an attack has been resolved, if the attacker deals damage to the target, the target first subtracts their \Block{}\index{Attribute!Block} from the damage rolled by the attacker and then any remaining damage is subtracted from its \Life{}\index{Attribute!Life}.

\subsection{Healing}\index{Healing}
To recover \Life{} lost during combat or otherwise during their adventures, creatures will need to \textbf{Heal}.
When a creature successfully heals, it regains a number of \Life{} equal to its \Recovery{}.
The easiest and simplest way to heal, is to get a night's rest; if a creature rests for 8 hours, with at least 6 hours of sleep and no strenuous activity over the period of time, they heal.
Creatures can attempt to heal, or attempt to help another creature heal, in less time with a roll, with the difficulty depending on the time taken to heal.
The \Difficulty{} is generally set based on the time taken to rest: if they spend 4 hours the roll is \emph{Very Easy}, 1 hour is \emph{Easy}, 30 minutes is \emph{Medium}, 10 minutes is \emph{Hard}, and 1 minute is \emph{Very Hard}.
The GM may adjust these difficulties based on the situation, for instance a hospital might decrease the difficulty while a busy or dangerous environment may increase it.
The time it takes to heal accounts both for time to tend wounds, relax, and so on, it also includes the time to find the necessary supplies; the GM may allow a shorter rest time without increasing the difficulty if sufficient supplies are available.
Based on the situation, the GM may allow healing to be attempted as an \emph{Action} in combat, though such a roll is likely to be \emph{Very Hard} and \emph{Decreased}.
On a healing roll, the target of the healing regains \Life{} equal to their \Recovery{} if the roll is \Good{} or \VeryGood{}, on a \Mixed{} result they heal half of their \Recovery{}.
On a \Bad{} or \VeryBad{} result, the target does not heal but on a \Bad{} result the GM should not impose a negative outcome, reserving that for a \VeryBad{} result.

A creature can only benefit from the effects of a night's rest or a successful roll for healing once per day each.

\subsubsection{Triaging}\index{Triaging}\label{sec:triage}
Sometime quicker healing is needed, especially if someone needs to be healed in the midst of combat.
To attempt healing in the midst of combat, the healer must spend an \Advantage{} and requires a relevant \Specialization{} that can be used to help heal.
The difficulty of the healing is generally \emph{Medium}; on a success the target gains number of \Life{} equal to the healer's specialization bonus.
On a success, the healer can also spend \Advantage{} to heal an additional \Life{} per \Advantage{}.

Once a creature is successfully healed in this manner attempts to heal them in this way again require spending an additional \Advantage.
If the \Specialization{} used has the Simple \& Weak tag, then the ability to use it for Triaging must be specifically defined as part of the ability.

\subsection{Defeat}
When a creature is reduced to 0 \Life{}, they are \textbf{Defeated}\index{Defeat}.
Depending on the circumstances, the GM may declare that the creature falls unconscious or becomes incapacitated for a number of hours, at which point they will awake with 1 \Life{}.
If the creature is still in danger, for instance in combat, the creature must choose one of the following options:
\begin{itemize}
  \item[-] \textbf{Unconsciousness}: the creature falls unconscious and is stable for the moment.
    An unconscious creature will remain unconscious for several hours, if not attended to, and then awaken with 1 \Life{}.
    If an unconscious creature is hit by an attack, it becomes \emph{Unstable}.
  \item[-] \textbf{Unstable}: the creature remains alive but barely, with its fate hanging in the balance.
    The creature remains conscious and can act.
    At the end of each of its turns and each time an unstable creature takes damage, it must roll using its \emph{Willpower} \Attribute{} and must achieve a \Mixed{} or better outcome; the \Difficulty{} of this role is \emph{Very Easy} the first time and increases one step each time they roll.
    If an unstable creature fails this roll, it dies.
    An \emph{Unstable} creature can be stabilized into \emph{Unconsciousness} by a successful \emph{Medium} \Difficulty{} roll.
\end{itemize}

\subsubsection{Optional Rule: Blaze of Glory}
A group may chose to play with the following third option for when a player character is defeated:
\begin{itemize}
  \item[-] \textbf{Blaze of Glory}: the creature dies, but gets to take a final Action whose result is \VeryGood.
    This action may not be used to Heal itself.
\end{itemize}

\subsubsection{Optional Rule: Resurrection}\index{Resurrection}
For groups that want it, death does not need to be the end of a character's journey, because there's magic in this world.
A group should discuss whether they want to include resurrection in their game before beginning a campaign and exactly how it should work, but the following rules are provided as a suggestion.

Each character gains three boxes, that are initially unmarked, and marking these boxes represents the character's connection to the world slipping away.
When a character dies, their party can attempt to resurrect them as an act of Big Magic, the magic should be made difficult not impossible, a Magnitude double the number of other party members is likely reasonable but the group should agree from the beginning.
For each roll to perform the Big Magic that is a result of \Bad{} or \VeryBad{}, the dead character marks one of the boxes.
If the character marks their third box they are permanently dead and cannot be restored to life as they have lost their connection to the world.
If a character is successfully resurrected, they regain 1 \Life{}.
The process of dying and being resurrected has an impact on a character; they take a permanent $-1$ penalty to one of their experiences and gain a $+1$ bonus (or new $+ 1$ experience).
They take an additional $-1$ penalty and $+1$ bonus for each box that was checked during their resurrection.
A group and GM may choose to allow a player to spend 3 XP to unmark a marked box, in which case the character also loses the associated $-1$ penalty, thoughts keeps the $+1$ bonus they received.

If using the optional Blaze of Glory option described above, it is recommended that that option also marks all of the boxes immediately, heroic sacrifice is a lot less heroic and meaningful when they can be brought back.

If a character marks their last box during an attempt to resurrect them, if the player is interested and with the agreement of the GM and the group, the player may create a new character that possesses the body of their former character, as the failed resurrection actually brought some other being into the body.
The new character should be created and leveled as normal, the player can choose whether or not to unmark the boxes and for each box left marked they retain a total of $+ 2$ experiences from the dead creature, representing how their new character still has some memory of the body's former occupant.
